\section{Interpretation}
\label{sec:interpretation}

The method can be viewed as a composite metric-space generator. The continuous
coordinates contribute standardised Euclidean directions, while each
categorical coordinate contributes a supervised block whose neighbourhoods are
optimised for predicting that categorical label from the remaining columns.
The final distance is therefore not a universal dataframe metric. It is an
empirical metric assembled from local prediction problems, and its validity
must be judged through the diagnostics in Section~\ref{sec:evaluation-framework}.

From a metric-learning perspective, the per-column NCA blocks are useful only
when they improve held-out categorical prediction relative to simpler
references. This is why the paper includes NCA-vs-PCA and majority-class
diagnostics. If a block improves prediction, then it provides evidence that the
supervised latent coordinates encode information relevant to that categorical
variable. If it does not, the block remains an implementation component but no
longer supports a mechanism claim.

From a manifold-augmentation perspective, the transport operator approximates
local tangent movement without explicitly estimating tangent planes or
curvature. If the fitted latent space is locally linear around the sampled
neighbour chain, then the displacement $C-B$ can be interpreted as a local
direction that may also be plausible near $A$. This assumption is weaker and
less formal than a tangent-space model, but it has an operational advantage:
the sampled anchor, neighbour chain, displacement, and decoded categorical
probabilities are all inspectable.

\subsection{Analogy-Based Interpretation}

The generation operator is not only a geometric mechanism; it is also the
primary explanation interface for individual samples. A generated row is
constructed from four concrete latent rows: the anchor $A$, the first observed
neighbour $B$, the second observed neighbour $C$, and the candidate
$D=T(A;B,C)$. The relation can be read as
\[
  D : A \approx C : B,
\]
meaning that $D$ changes from $A$ in the same fitted-latent direction that
$C$ changes from $B$. This is a data-driven displacement analogy. It does not
assert that the change is causal, semantically named, or valid outside the
local metric context; it only records how one observed local difference was
reused to construct a new row.

\begin{table}[t]
\centering
\caption{Analogy trace for one generated row. The entries are roles in the
fitted latent construction, not domain-level causal explanations. Takeaway:
individual samples are inspectable because each one can be traced to a concrete
anchor and an observed neighbour-pair displacement.}
\label{tab:analogy-trace}
\small
\begin{adjustbox}{max width=\textwidth}
\begin{tabular}{@{}lll@{}}
\toprule
Symbol & Role & Interpretation \\
\midrule
$A$ & Anchor row & Observed row to be modified \\
$B$ & Source neighbour & Observed row where the local displacement starts \\
$C$ & Second neighbour & Observed row where the local displacement ends \\
$C-B$ & Displacement & Fitted-latent change transferred to the anchor \\
$D=T(A;B,C)$ & Generated candidate & Anchor modified by the observed displacement \\
\bottomrule
\end{tabular}
\end{adjustbox}
\end{table}

This analogy trace distinguishes the sampler from generators that expose only
global fitted parameters or a sampled latent vector. The trace is not enough to
prove realism, but it gives a user a concrete object to inspect when a row looks
implausible: the issue can be localised to the anchor choice, the neighbour
pair, the scale $\lambda$, the rejection constraints, or the decoder
probabilities.

Table~\ref{tab:adult-analogy-example} shows one concrete Adult Census Income
trace from a deterministic fitted run. The row values are displayed only for a
small set of human-readable columns; the actual operator is applied in the
fitted latent space before decoding. This example illustrates the analogy
$D : A \approx C : B$ in concrete tabular terms.

\begin{table}[t]
\centering
\caption{Concrete Adult analogy trace for one generated row. The neighbour
pair changes from a 50-year-old, high-income handlers-cleaners row to a
56-year-old, lower-income craft-repair row; the generated row applies a similar
observed displacement to the 62-year-old anchor. Takeaway: the generation step
can be read as a traceable observed-displacement analogy over real examples.}
\label{tab:adult-analogy-example}
\small
\begin{adjustbox}{max width=\textwidth}
\begin{tabular}{@{}lp{0.72\textwidth}@{}}
\toprule
Role & Displayed values \\
\midrule
$A$ (anchor) & age=62, education=HS-grad, occupation=Farming-fishing, hours/week=40, income=$>$50K \\
$B$ (neighbour) & age=50, education=HS-grad, occupation=Handlers-cleaners, hours/week=40, income=$>$50K \\
$C$ (neighbour) & age=56, education=HS-grad, occupation=Craft-repair, hours/week=40, income=$\leq$50K \\
$D$ (generated) & age=68, education=HS-grad, occupation=Craft-repair, hours/week=40, income=$\leq$50K \\
\bottomrule
\end{tabular}
\end{adjustbox}
\end{table}

\subsection{Relation to Local Linear Transport}

The operator $T(A;B,C)=A+\lambda(C-B)$ can be read as applying a learned local
displacement field in the fitted metric space. The pair $(B,C)$ supplies an
observed first-order direction in one neighbourhood, and the anchor $A$ is the
point to which that direction is transported. Under a local linearity
assumption, nearby displacement vectors approximate directions tangent to the
data support. DataFrameSampler does not estimate a tangent basis or curvature;
instead it samples one observed displacement and reuses it. The resulting
operator is therefore a structured approximation to local transport, not a
claim that the fitted space is globally flat.

This view also clarifies the role of $\lambda$. Small values keep candidates
closer to the anchor and behave more like conservative neighbourhood
perturbations. Larger values transfer more of the neighbour-chain displacement
and can produce out-of-hull candidates, but also increase the chance of
leaving the useful local region. The rejection constraints are a practical
guardrail around this operator; they are not a projection onto the true data
manifold.

\subsection{Implicit Objective View}

Although the method has no single global likelihood, adversarial loss, or
energy function, it has a coherent proxy interpretation. The NCA blocks seek
local label consistency for each categorical column under a supervised metric.
The decoder then asks whether the final context still predicts the held-out
categorical label without seeing the label's own block. The transport step
samples displacements that preserve observed local neighbour-chain statistics
in that same fitted metric space. In this sense, the pipeline approximates the
preservation of local neighbourhood structure under a composite supervised
metric. The mechanism evidence in the experiments tests each part of this
proxy interpretation rather than treating it as an optimised objective.

Two lightweight sanity expectations follow from this interpretation. First,
when NCA improves class separation for a categorical column, the corresponding
decoder should in many cases face a lower-ambiguity prediction problem,
although this is an empirical expectation rather than a theorem. Second, when
the latent space is locally regular, out-of-hull transported points should be
insertable into a frozen manifold geometry with stress comparable to held-out
real points. The paper tests both expectations diagnostically; it does not use
them as guarantees.
